\section{Soundness}
\label{sec:soundness}

This section proves the main results: every axiom of $\ThE$ holds in
$\ML$, and consequently first-order provability of $\F{G}$ entails truth
of $G$ in the semantics of \cref{sec:semantics}.

\subsection{Guard disciplines}
\label{sec:soundness-discipline}

\begin{definition}[Context satisfaction]
\label{def:ctx-sat}
A valuation $\rho$ \emph{satisfies} a context $\Gamma$, written $\rho
\vDash \Gamma$, if in telescope order: $\rho(x) \in \semr{\sigma}{\rho}$
for $x : \sigma$; $\rho(X) : \Dom^n \to \{0,1\}$ for $X :
(\tup{\sigma}) \to \Prop$; and $\pvr{\varphi}{\rho} = 1$ for $p :
\varphi$.
\end{definition}

The treatment of $\tfix$ interpolates between ``$\rho(f)$ is anything''
and ``$\rho(f)$ inhabits its type'' by a rank-restricted membership; the
only change against \cite{Extraction2026} is that the rank of the guard
argument is its \emph{element rank} in the guard family
(\cref{def:rank}), which \cref{lem:index-reflection} makes well defined
without reference to the indices at which the argument is used.

\begin{lemma}[Flattening of telescopes]
\label{lem:flatten}
Let $T = \Pi x_1{:}\sigma_1 \dots \Pi x_n{:}\sigma_n.\tau$.  Then $F \in
\semr{T}{\rho}$ iff for all $d_1, \dots, d_n$ chosen successively with
$d_i \in \semr{\sigma_i}{\rho[x_1 \mapsto d_1, \dots, x_{i-1} \mapsto
d_{i-1}]}$ we have $F \cdot d_1 \cdots d_n \in
\semr{\tau}{\rho[\tup{x} \mapsto \tup{d}\,]}$.
\end{lemma}

\begin{proof}
Unfold \cref{def:interp} $n$ times.
\end{proof}

\begin{definition}[Rank-restricted interpretation]
\label{def:rank-restricted}
Let $T = \Pi x_1{:}\sigma_1 \dots \Pi x_n{:}\sigma_n.\tau$ be a fix
type with guard position $k$, $\sigma_k = I(\tup{v})$, and let $m \in
\mathbb{N}$.  Define $\semr{T}{\rho}^{k, <m}$ as the set of $F \in
\Dom$ such that for all successively chosen $d_1, \dots, d_n$ as in
\cref{lem:flatten} with additionally $\rank{I}{d_k} < m$:
\[
  F \cdot d_1 \cdots d_n \in \semr{\tau}{\rho[\tup{x} \mapsto
  \tup{d}\,]}.
\]
(The rank is defined: the chain condition places $d_k$ in
$\semr{I(\tup{v})}{\rho[\dots]}$, so $(\cdot, d_k) \in \semr{I}{}$.)  By
\cref{lem:flatten} and finiteness of ranks, $\semr{T}{\rho} =
\bigcap_{m} \semr{T}{\rho}^{k,<m}$, and $\semr{T}{\rho}^{k,<m}$ is
antitone in $m$.
\end{definition}

\begin{definition}[Guard discipline]
\label{def:discipline}
A \emph{guard discipline} $\mathcal{D}$ over $\Gamma$ is a finite set of
triples $(f, V_f, m_f)$ where the $f$ are pairwise distinct term
variables declared in $\Gamma$ with fix types $T_f$ (guard position
$k_f$, guard family $I_f$), $V_f$ is a set of term variables declared in
$\Gamma$ whose types are occurrences of $I_f$, and $m_f \in \mathbb{N}$.
A phrase $t$ \emph{respects} $\mathcal{D}$ if for each $(f, V_f, m_f)
\in \mathcal{D}$, every occurrence of $f$ in $t$ lies at the head of an
application spine $f\,a_1 \cdots a_{n_f}$ occurring as a subterm of $t$,
with $a_{k_f} \in V_f$.  A valuation $\rho$ satisfies $\Gamma$
\emph{under} $\mathcal{D}$, written $\rho \vDash_{\mathcal{D}} \Gamma$,
if $\rho$ satisfies all declarations of $\Gamma$ as in
\cref{def:ctx-sat} except that for each $(f, V_f, m_f) \in \mathcal{D}$
only $\rho(f) \in \semr{T_f}{\rho}^{k_f, < m_f}$ is required, and
additionally $\rank{I_f}{\rho(y)} < m_f$ for every $y \in V_f$.
\end{definition}

\begin{lemma}[Locality of respect]
\label{lem:respect-local}
Let $t$ respect $\mathcal{D}$ and let $s$ be a subterm of $t$ occurring
in a non-applied position (any position other than the function part of
an application).  Then $s$ respects $\mathcal{D}$.
\end{lemma}

\begin{proof}
As in \cite{Extraction2026}: if the witnessing spine of an occurrence of
$f$ in $s$ were not contained in $s$, then $s$ would be a proper prefix
of that spine and hence occur applied, contradicting the assumption.
\end{proof}

\subsection{The fundamental lemma}
\label{sec:soundness-fundamental}

Throughout this subsection we assume the \emph{environment hypothesis}
for the fixed environment $\Env$ (discharged in
\cref{thm:env-validity}):
\begin{itemize}[leftmargin=2em]
\item[(H)] for every term definition $c \defeq t : \tau_c$ in $\Env$,
  $\cls{c} \in \semr{\tau_c}{}$; and for every lemma $c \defeq \pi :
  \varphi_c$ in $\Env$, $\pv{\varphi_c} = 1$.
\end{itemize}

We use the simultaneous form of \cref{lem:sem-subst} throughout: for a
simultaneous substitution of terms $\tup{a}$ for distinct variables
$\tup{x}$, $\semr{\sigma[\tup{a}/\tup{x}]}{\rho} =
\semr{\sigma}{\rho[\tup{x} \mapsto \sem{\tup{a}}_\rho]}$, obtained from
the single-variable form by $\alpha$-renaming to fresh variables and
iterating; likewise for terms and propositions.

\begin{lemma}[Fundamental lemma]
\label{lem:fundamental}
Assume (H).  Let $\mathcal{D}$ be a guard discipline over $\Gamma$ and
$\rho \vDash_{\mathcal{D}} \Gamma$.
\begin{enumerate}[leftmargin=2em]
\item If $\Gamma \vdash t : \sigma$ and $t$ respects $\mathcal{D}$, then
  $\sem{t}_\rho \in \semr{\sigma}{\rho}$.
\item If $\Gamma \vdash \pi : \varphi$ and $\pi$ respects $\mathcal{D}$,
  then $\pvr{\varphi}{\rho} = 1$.
\end{enumerate}
\end{lemma}

\begin{proof}
By induction on the size of the typing derivation, mutually for (1) and
(2).  We freely use \cref{lem:sem-basic} (weakening),
\cref{lem:sem-subst} (semantic substitution, in simultaneous form),
\cref{lem:conv-inv} (conversion invariance, which handles the two
conversion rules), and \cref{lem:respect-local} to propagate the
discipline to subphrases; types, propositions and predicate expressions
contain no applied occurrences of disciplined variables in well-formed
judgements, by the same locality argument.  When extending $\rho$ at
bound variables we $\alpha$-rename so that disciplined variables,
$V_f$-variables and context variables are not shadowed.

\emph{Variables, constants.}  If $t = x$ is not disciplined, $\rho(x)
\in \semr{\sigma}{\rho}$ by $\rho \vDash_{\mathcal{D}} \Gamma$; a bare
disciplined $f$ does not respect $\mathcal{D}$, so this case cannot
arise for it.  For a term constant, $\sem{c}_\rho = \cls{c} \in
\semr{\tau_c}{}$ by (H), the $\delta$-rule and weakening; for a lemma
constant in (2), directly by (H).

\emph{Disciplined spines.}  Suppose $t = f\,a_1 \cdots a_j$ with $(f,
V_f, m_f) \in \mathcal{D}$ and $j \geq n_f$ (respect makes $f\,a_1
\cdots a_{n_f}$ a spine of $t$), and let $T_f = \Pi
x_1{:}\sigma_1 \dots \Pi x_{n_f}{:}\sigma_{n_f}.\tau$.  The typing
derivation decomposes into $j$ applications; by induction hypothesis
(1) applied to each argument (which respects $\mathcal{D}$ by
\cref{lem:respect-local}), $d_i \defeq \sem{a_i}_\rho \in
\semr{\sigma_i}{\rho[x_1 \mapsto d_1, \dots]}$ for $i \le n_f$, the
intermediate assignments being handled by semantic substitution as in
the application case below.  By respect, $a_{k_f} \in V_f$, so
$\rank{I_f}{d_{k_f}} = \rank{I_f}{\rho(a_{k_f})} < m_f$; since
$\rho(f) \in \semr{T_f}{\rho}^{k_f, <m_f}$,
\[
  \sem{f\,a_1 \cdots a_{n_f}}_\rho = \rho(f) \cdot d_1 \cdots d_{n_f}
  \in \semr{\tau}{\rho[\tup{x} \mapsto \tup{d}\,]}
  = \semr{\subst{\tau}{\tup{a}}{\tup{x}}}{\rho},
\]
the last step by semantic substitution; remaining arguments $a_{n_f+1},
\dots, a_j$ are absorbed by the application case.

\emph{Applications and abstractions.}  For $t\,a$ with head not
disciplined: the function part respects $\mathcal{D}$ (its disciplined
occurrences have their spines within it, by the argument of
\cref{lem:respect-local} applied to the maximal spine decomposition),
so by the induction hypotheses $\sem{t}_\rho \in \semr{\Pi
x{:}\sigma.\tau}{\rho}$ and $\sem{a}_\rho \in \semr{\sigma}{\rho}$,
whence $\sem{t\,a}_\rho = \sem{t}_\rho \cdot \sem{a}_\rho \in
\semr{\tau}{\rho[x \mapsto \sem{a}_\rho]} =
\semr{\subst{\tau}{a}{x}}{\rho}$.  For $t\,\pi$: $\erase{t\,\pi} =
\erase{t}$; induction hypothesis (2) gives $\pvr{\varphi}{\rho} = 1$,
so induction hypothesis (1) and the clause for $\semr{\varphi \to
\tau}{\rho}$ conclude.  For $\lambda x{:}\sigma.t$ and $e \in
\semr{\sigma}{\rho}$: $\sem{\lambda x{:}\sigma.t}_\rho \cdot e =
\cls{(\lambda x.\erase{t}^\rho)\,\hat{e}} = \sem{t}_{\rho[x \mapsto
e]}$ (one $\beta$-step), which lies in $\semr{\tau}{\rho[x \mapsto e]}$
by the induction hypothesis.  For $\lambda p{:}\varphi.t$: if
$\pvr{\varphi}{\rho} = 1$ then $\rho \vDash_{\mathcal{D}} \Gamma,
p{:}\varphi$ and $\sem{\lambda p{:}\varphi.t}_\rho = \sem{t}_\rho \in
\semr{\tau}{\rho}$ by the induction hypothesis, as the clause for
$\varphi \to \tau$ requires.

\emph{Constructors.}  Let $t = \wcns{c}_i(\tup{a}; \tup{\pi})$ with the
typing premises of \cref{fig:terms}, so the displayed type is
$I(\tup{t}_i[\tup{a}/\tup{y}])$.  Let $\rho_0 \defeq [\tup{y} \mapsto
\sem{\tup{a}}_\rho]$.  By induction hypothesis (1) on each $a_j$ and
semantic substitution, $\sem{a_j}_\rho \in
\semr{A_{ij}[\tup{a}_{<j}/\tup{y}_{<j}]}{\rho} =
\semr{A_{ij}}{\rho_0}$; for recursive $j$ with $A_{ij} =
I(\tup{v}_{ij})$ this reads $(\sem{\tup{v}_{ij}}_{\rho_0},
\sem{a_j}_\rho) \in \semr{I}{}$.  By induction hypothesis (2) on each
$\pi_l$ and semantic substitution,
$\pvr{\varphi_{il}}{\rho_0} =
\pvr{\varphi_{il}[\tup{a}/\tup{y}]}{\rho} = 1$.  By
\cref{lem:index-reflection}(5),
\[
  \bigl(\sem{\tup{t}_i}_{\rho_0},\, \cls{\wcns{c}_i(\tup{e})}\bigr) \in
  \semr{I}{},
  \qquad \text{where } \cls{\wcns{c}_i(\tup{e})} = \sem{t}_\rho,
\]
and $\sem{\tup{t}_i}_{\rho_0} =
\sem{\tup{t}_i[\tup{a}/\tup{y}]}_{\rho}$ by semantic substitution, so
$\sem{t}_\rho \in \semr{I(\tup{t}_i[\tup{a}/\tup{y}])}{\rho}$.

\emph{Forded case.}  Let $t = \tcase{a}{\tup{z}.z.\tau}
{\{\wcns{c}_i(\tup{y}_i;\tup{p}_i;\tup{q}_i) \Rightarrow t_i\}_i}$ with
$\Gamma \vdash a : I(\tup{u})$, and displayed type $\tau[\tup{u}/\tup{z},
a/z]$.  By induction hypothesis (1), $(\sem{\tup{u}}_\rho, \sem{a}_\rho)
\in \semr{I}{}$.  \Cref{lem:index-reflection}(1) yields the unique
decomposition: an $i$, classes $\cls{\tup{e}}$ with $\sem{a}_\rho =
\cls{\wcns{c}_i(\tup{e})}$, the memberships and premise truths at
$\rho_0 = [\tup{y}_i \mapsto \cls{\tup{e}}\,]$, and
\begin{equation}
  \sem{\tup{u}}_\rho = \sem{\tup{t}_i}_{\rho_0}.
  \tag{$*$}
\end{equation}
Let $\rho' \defeq \rho[\tup{y}_i \mapsto \cls{\tup{e}}\,]$ ($\tup{y}_i$
fresh).  Then $\rho' \vDash_{\mathcal{D}} \Gamma, \tup{y}_i{:}\tup{A}_i,
\tup{p}_i{:}\tup{\varphi}_i, \tup{q}_i{:}(\tup{u} =_{\tup{\sigma}}
\tup{t}_i)$: the $\Gamma$-part by weakening; the memberships and
premises by the decomposition (their free variables are among earlier
declarations and $\tup{y}_i$, on which $\rho'$ agrees with $\rho_0$);
and for the index equations,
$\pvr{u_h =_{\sigma_h} t_{ih}}{\rho'} = 1$ because $\sem{u_h}_{\rho'} =
\sem{u_h}_\rho$ ($\tup{y}_i$ fresh for $u_h$) $=
\sem{t_{ih}}_{\rho_0} = \sem{t_{ih}}_{\rho'}$ by ($*$).  On erasures,
\[
  \erase{t}^{\rho} = \mathsf{case}(\erase{a}^{\rho}; \dots)
  \lconv \mathsf{case}(\wcns{c}_i(\tup{e}); \dots)
  \lred \erase{t_i}^{\rho'},
\]
so $\sem{t}_\rho = \sem{t_i}_{\rho'}$.  By the induction hypothesis on
the branch derivation ($t_i$ respects $\mathcal{D}$ by
\cref{lem:respect-local}),
\[
  \sem{t_i}_{\rho'} \in
  \semr{\tau[\tup{t}_i/\tup{z},\,
    \wcns{c}_i(\tup{y}_i;\tup{p}_i)/z]}{\rho'}
  = \semr{\tau}{\rho'[\tup{z} \mapsto \sem{\tup{t}_i}_{\rho'},\,
      z \mapsto \sem{\wcns{c}_i(\tup{y}_i;\tup{p}_i)}_{\rho'}]}
\]
by semantic substitution.  Now $\sem{\tup{t}_i}_{\rho'} =
\sem{\tup{t}_i}_{\rho_0} = \sem{\tup{u}}_\rho$ by ($*$), and
$\sem{\wcns{c}_i(\tup{y}_i;\tup{p}_i)}_{\rho'} =
\cls{\wcns{c}_i(\tup{e})} = \sem{a}_\rho$; by weakening on the remaining
free variables of $\tau$ and semantic substitution backwards, the
displayed set equals $\semr{\tau[\tup{u}/\tup{z}, a/z]}{\rho}$, as
required.  This is the place where the index equations of the model ---
\cref{lem:index-reflection}(1) --- convert the branch's type, stated at
the constructor's indices, to the case's type, stated at the
occurrence's indices; in CIC the same conversion is performed by index
unification, and in the forded reading by rewriting along $\tup{q}_i$.

\emph{Fix.}  Let
$t = \tfix\,f{:}T.\,\lambda \tup{x}.\,\tcase{x_k}{\tup{z}.z.\tau_0}
{\{\wcns{c}_i(\tup{y}_i;\tup{p}_i;\tup{q}_i) \Rightarrow b_i\}_i}$ with
$T = \Pi x_1{:}\sigma_1 \dots \Pi x_n{:}\sigma_n.\tau$, $\sigma_k =
I(\tup{v})$, $\tau = \tau_0[\tup{v}/\tup{z}, x_k/z]$, and let $F \defeq
\sem{t}_\rho$.  Since $\semr{T}{\rho} = \bigcap_m
\semr{T}{\rho}^{k,<m}$, it suffices to prove $F \in
\semr{T}{\rho}^{k,<m}$ for all $m$, by strong induction on $m$ (the
\emph{rank induction}, nested inside the main induction).

Let $m$ be given with $F \in \semr{T}{\rho}^{k,<m'}$ for all $m' < m$,
and let $d_1, \dots, d_n$ be a chain as in \cref{def:rank-restricted}
with $m_0 \defeq \rank{I}{d_k} < m$.  Write $\rho_c \defeq \rho[x_1
\mapsto d_1, \dots, x_{k-1} \mapsto d_{k-1}]$; the chain condition gives
$(\sem{\tup{v}}_{\rho_c}, d_k) \in \semr{I}{}$.  By
\cref{lem:index-reflection}(1) there are $i$ and $\cls{\tup{e}}$ with
$d_k = \cls{\wcns{c}_i(\tup{e})}$, the memberships and premises at
$\rho_0 = [\tup{y}_i \mapsto \cls{\tup{e}}\,]$, and
$\sem{\tup{v}}_{\rho_c} = \sem{\tup{t}_i}_{\rho_0}$; by
\cref{lem:index-reflection}(4), $\rank{I}{\cls{e_j}} < m_0$ for
recursive $j$.  Unfolding one $\mu$-step, $n$ $\beta$-steps and one
$\iota$-step on a representative of $F \cdot d_1 \cdots d_n$:
\[
  F \cdot d_1 \cdots d_n = \sem{b_i}_{\rho''}, \qquad
  \rho'' \defeq \rho[f \mapsto F,\ \tup{x} \mapsto \tup{d},\
  \tup{y}_i \mapsto \cls{\tup{e}}\,].
\]
Apply the \emph{main} induction hypothesis to the branch derivation
\[
  \Gamma, f{:}T, \tup{x}{:}\tup{\sigma},
  \tup{y}_i{:}\tup{A}_i, \tup{p}_i{:}\tup{\varphi}_i,
  \tup{q}_i{:}(\tup{v} =_{\tup{\sigma}'} \tup{t}_i)
  \vdash b_i : \tau_0[\tup{t}_i/\tup{z},\,
    \wcns{c}_i(\tup{y}_i;\tup{p}_i)/z]
\]
with discipline $\mathcal{D}' = \mathcal{D} \cup \{(f,
\tup{y}_i^{\mathrm{rec}}, m_0)\}$ and valuation $\rho''$.  The premises
hold: $b_i$ respects $\mathcal{D}'$ (for $f$ by the guard condition
$\mathrm{grd}_k$; for the outer disciplines by
\cref{lem:respect-local}); $\rho''(f) = F \in \semr{T}{\rho}^{k,<m_0}$
by the rank induction hypothesis; $\rank{I}{\rho''(y)} < m_0$ for $y \in
\tup{y}_i^{\mathrm{rec}}$ as noted; the $\tup{x}$- and
$\tup{y}_i$-declarations are satisfied by the chain condition and the
decomposition; the premise proofs $\tup{p}_i$ by the decomposition; and
the index equations $\tup{q}_i$ by $\sem{\tup{v}}_{\rho''} =
\sem{\tup{v}}_{\rho_c} = \sem{\tup{t}_i}_{\rho_0} =
\sem{\tup{t}_i}_{\rho''}$ ($\tup{v}$ mentions only $x_1, \dots,
x_{k-1}$).  The conclusion gives $\sem{b_i}_{\rho''} \in
\semr{\tau_0[\tup{t}_i/\tup{z}, \wcns{c}_i(\tup{y}_i;\tup{p}_i)/z]}
{\rho''}$; by semantic substitution this set is
$\semr{\tau_0}{\rho''[\tup{z} \mapsto \sem{\tup{t}_i}_{\rho''}, z
\mapsto d_k]} = \semr{\tau_0}{\rho''[\tup{z} \mapsto
\sem{\tup{v}}_{\rho''}, z \mapsto \sem{x_k}_{\rho''}]} =
\semr{\tau_0[\tup{v}/\tup{z}, x_k/z]}{\rho''} =
\semr{\tau}{\rho[\tup{x} \mapsto \tup{d}\,]}$ (weakening: $f$,
$\tup{y}_i$ do not occur in $\tau$).  This closes the rank induction,
hence the fix case.

\emph{Refinement formers.}  $\trefine{t}{\pi}$: by the induction
hypotheses, $\sem{t}_\rho \in \semr{\sigma}{\rho}$ and
$\pvr{\subst{\varphi}{t}{x}}{\rho} = 1$, i.e.\
$\pvr{\varphi}{\rho[x \mapsto \sem{t}_\rho]} = 1$ by semantic
substitution; since $\erase{\trefine{t}{\pi}} = \erase{t}$, this is
membership in $\semr{\refty{x{:}\sigma}{\varphi}}{\rho}$.  $\tval{t}$:
membership in the refinement interpretation includes membership in
$\semr{\sigma}{\rho}$, and $\erase{\tval{t}} = \erase{t}$.  For
$\tvalid{t}$ in (2): the payload component of the same membership,
rewritten by semantic substitution using $\sem{\tval{t}}_\rho =
\sem{t}_\rho$.

\emph{Propositional singleton eliminations.}  $\tsplit{\pi}{p_1
p_2.t}$: induction hypothesis (2) gives both conjuncts, so $\rho
\vDash_{\mathcal{D}} \Gamma, p_1{:}\varphi_1, p_2{:}\varphi_2$;
conclude by induction hypothesis (1) on $t$, the erasures agreeing.
$\tcast{\pi}{z.\tau}{t}$: induction hypothesis (2) gives $\sem{a}_\rho
= \sem{b}_\rho$, hence $\semr{\subst{\tau}{a}{z}}{\rho} =
\semr{\tau}{\rho[z \mapsto \sem{a}_\rho]} =
\semr{\subst{\tau}{b}{z}}{\rho}$; conclude by induction hypothesis (1)
on $t$.  $\texfalso{\pi}{\sigma}$ (and $\texfalso{\pi}{\varphi}$ in
(2)): vacuous --- induction hypothesis (2) would give
$\pvr{\bot}{\rho} = 1$.

\emph{Singleton elimination.}  Let $t =
\tscase{J}{\pi}{\tup{z}.\tau}{\tup{p}.t_0}$ with $J$ a forced singleton
(clause $\wcns{k} : \forall \tup{x}{:}\tup{\tau}.\ \varphi_1 \to \dots
\to \varphi_\ell \to J(\tup{u}^0)$, solution $\jmath$), $\Gamma \vdash
\pi : J(\tup{u})$, $\theta = [u_{\jmath(i)}/x_i]_i$, and displayed type
$\tau[\tup{u}/\tup{z}]$.  By induction hypothesis (2),
$\sem{\tup{u}}_\rho \in \pv{J}$; unfolding the least stage containing
the tuple (the predicate analogue of
\cref{lem:index-reflection}(1), immediate from \cref{def:interp} since
$J$ has a single clause): there is $\rho_0 = [\tup{x} \mapsto \tup{d}\,]$
with $d_i \in \semr{\tau_i}{\rho_0}$, $\pvr{\varphi_l}{\rho_0} = 1$ for
every $l$ (for a recursive premise $J(\tup{v})$, membership of
$\sem{\tup{v}}_{\rho_0}$ in an earlier stage gives exactly
$\pvr{J(\tup{v})}{\rho_0} = 1$), and $\sem{\tup{u}}_\rho =
\sem{\tup{u}^0}_{\rho_0}$.  By forcedness, for each $i$:
\[
  \sem{\theta(x_i)}_\rho = \sem{u_{\jmath(i)}}_\rho =
  \bigl(\sem{\tup{u}^0}_{\rho_0}\bigr)_{\jmath(i)} =
  \sem{u^0_{\jmath(i)}}_{\rho_0} = \sem{x_i}_{\rho_0} = d_i .
\]
Hence by semantic substitution $\pvr{\varphi_l[\theta]}{\rho} =
\pvr{\varphi_l}{\rho[\tup{x} \mapsto \tup{d}\,]} =
\pvr{\varphi_l}{\rho_0} = 1$, so $\rho \vDash_{\mathcal{D}} \Gamma,
\tup{p}{:}\theta(\tup{\varphi})$, and the induction hypothesis (1) on
$t_0$ gives $\sem{t_0}_\rho \in
\semr{\tau[\theta(\tup{u}^0)/\tup{z}]}{\rho} = \semr{\tau}{\rho[\tup{z}
\mapsto \sem{\theta(\tup{u}^0)}_\rho]}$.  Again by the displayed
computation, $\sem{\theta(u^0_h)}_\rho = \sem{u^0_h}_{\rho_0} =
\sem{u_h}_\rho$ for every $h$, so this set is
$\semr{\tau[\tup{u}/\tup{z}]}{\rho}$; and $\sem{t}_\rho =
\sem{t_0}_\rho$ since $\erase{t} = \erase{t_0}$.

\emph{Discrimination.}  For $\mathsf{discr}(\pi)$ with $\Gamma \vdash
\pi : \wcns{c}(\tup{a};\cdot) =_\sigma \wcns{d}(\tup{b};\cdot)$,
$\wcns{c} \neq \wcns{d}$: induction hypothesis (2) would give
$\cls{\wcns{c}(\erase{\tup{a}}^\rho)} =
\cls{\wcns{d}(\erase{\tup{b}}^\rho)}$, contradicting
\cref{cor:discrimination}(1); the case is vacuous --- no $\rho
\vDash_{\mathcal{D}} \Gamma$ reaches it.

\emph{Proof formers.}  The introduction and elimination rules for
$\to$, $\wedge$, $\vee$, $\forall$, $\exists$, $\top$ and predicate
quantification are sound for the truth-value semantics by direct
computation: each rule matches the corresponding $\min$/$\max$ clause
of \cref{def:interp}; the instantiation rules use semantic
substitution (parts (1) and (3) of \cref{lem:sem-subst}); and the
witnesses in $\mathsf{ecase}$ and predicate-quantifier elimination are
semantic (an element $d \in \semr{\sigma}{\rho}$, a function
$\kv{P}_\rho$), requiring no syntactic counterpart.
$\mathsf{refl}_\sigma(a)$ is reflexivity of class equality;
$\mathsf{rew}(\pi; z.\psi; \pi')$ uses $\sem{a}_\rho = \sem{b}_\rho$
and semantic substitution exactly as $\tcast{}{}{}$.

\emph{Predicate-constructor introduction.}  For
$\wcns{k}(\tup{a};\tup{\pi})$: by induction hypothesis (1) and iterated
semantic substitution along the dependent telescope, each
$\sem{a_i}_\rho$ lies in $\semr{\tau_i}{\rho_0}$ where $\rho_0 =
[\tup{x} \mapsto \sem{\tup{a}}_\rho]$; by induction hypothesis (2) the
non-recursive premises are true at $\rho_0$ and each recursive premise
$J(\tup{v})$ yields $\sem{\tup{v}}_{\rho_0} \in \pv{J}$, hence in some
finite stage.  The stage construction (at the maximum of those stages)
puts $\sem{\tup{u}}_{\rho_0}$ in $\pv{J}$, and semantic substitution
identifies it with $\sem{\subst{\tup{u}}{\tup{a}}{\tup{x}}}_\rho$.

\emph{Predicate induction.}  For $\tind{J}(\tup{z}.\psi; \tup{\pi};
\pi_0)$ with $\pi_0 : J(\tup{a})$: by induction hypothesis (2),
$\sem{\tup{a}}_\rho \in \pv{J}$ and $\pvr{\Xi_j(\tup{z}.\psi)}{\rho} =
1$ for every clause $j$.  We show by strong induction on $s$ that every
$\tup{d} \in \pv{J}^s$ satisfies $\pvr{\psi}{\rho[\tup{z} \mapsto
\tup{d}\,]} = 1$: a tuple enters $\pv{J}^{s}$ (with $s = s'+1$) through
some clause $j$ with witnesses $\tup{d}' \in
\semr{\tup{\tau}_j}{}$-telescope, true non-recursive premises, and
recursive-premise tuples in $\pv{J}^{s'}$; instantiating the semantic
universal quantifiers and implications of $\pvr{\Xi_j}{\rho} = 1$ at
these witnesses --- the inductive-hypothesis premises supplied by the
inner induction at $s'$ --- yields
$\pvr{\subst{\psi}{\tup{u}_j}{\tup{z}}}{\rho[\tup{x} \mapsto
\tup{d}']} = 1$, which by semantic substitution is
$\pvr{\psi}{\rho[\tup{z} \mapsto \tup{d}\,]} = 1$.  Conclude at
$\sem{\tup{a}}_\rho$.

\emph{Datatype induction.}  Let $\pi = \tind{I}(\tup{z}.z.\psi;
(\tup{y}_i\,\tup{p}_i\,\tup{h}_i.\pi_i)_i; a)$ with $\Gamma \vdash a :
I(\tup{u})$ and the clause derivations of \cref{fig:proofs2}.  By
induction hypothesis (1), $(\sem{\tup{u}}_\rho, \sem{a}_\rho) \in
\semr{I}{}$.  We show by strong induction on $\rank{I}{d}$ that every
$(\tup{d}, d) \in \semr{I}{}$ satisfies $\pvr{\psi}{\rho[\tup{z} \mapsto
\tup{d}, z \mapsto d]} = 1$.  Decompose $(\tup{d}, d)$ by
\cref{lem:index-reflection}(1) into $i$, $\cls{\tup{e}}$, $\rho_0$; for
each recursive $j$, $(\sem{\tup{v}_{ij}}_{\rho_0}, \cls{e_j}) \in
\semr{I}{}$ with $\rank{I}{\cls{e_j}} < \rank{I}{d}$
(\cref{lem:index-reflection}(4)), so the inner induction gives
$\pvr{\psi}{\rho[\tup{z} \mapsto \sem{\tup{v}_{ij}}_{\rho_0}, z \mapsto
\cls{e_j}]} = 1$, i.e., by semantic substitution and weakening,
$\pvr{\psi[\tup{v}_{ij}/\tup{z}, y_{ij}/z]}{\rho'} = 1$ at $\rho' \defeq
\rho[\tup{y}_i \mapsto \cls{\tup{e}}\,]$.  Together with the memberships
and premise truths of the decomposition, $\rho' \vDash_{\mathcal{D}}$
the clause context, so the main induction hypothesis (2) applied to
$\pi_i$ gives $\pvr{\psi[\tup{t}_i/\tup{z},
\wcns{c}_i(\tup{y}_i;\tup{p}_i)/z]}{\rho'} = 1$, which by semantic
substitution is $\pvr{\psi}{\rho[\tup{z} \mapsto \tup{d}, z \mapsto d]}
= 1$ (using $\sem{\tup{t}_i}_{\rho_0} = \tup{d}$ and
$\sem{\wcns{c}_i(\tup{y}_i;\tup{p}_i)}_{\rho'} = d$).  Conclude at
$(\sem{\tup{u}}_\rho, \sem{a}_\rho)$.
\end{proof}

\subsection{Validity of the emitted theory}
\label{sec:soundness-validity}

\begin{lemma}[Validity of specification axioms]
\label{lem:spec-valid}
Let $c \defeq t : \tau$ be a term definition with $\tau \in \Sfrag$, and
suppose $\cls{c} \in \semr{\tau}{}$.  Then $\ML \vDash \specf(c)$.
\end{lemma}

\begin{proof}
By induction on the type $\tau'$ in the recursion of \cref{def:spec},
with the invariant: if
$\rho$ assigns the accumulated variables, $v \defeq \cls{c\,\hat{e}_1
\cdots \hat{e}_j}$ for representatives of the accumulated values, and
$v \in \semr{\tau'}{\rho}$, then $\ML, \rho \vDash S(u; \tau')$.  The
$\Pi$-case extends $\rho$ under a guard, which \cref{lem:coincidence}
converts to a membership; the premise case uses \cref{lem:coincidence}
for $\F{\varphi}$ and leaves the accumulator unchanged; the base case
(datatype occurrence or refinement type) is \cref{lem:coincidence}(1)
--- for a family occurrence this includes the index arguments of the
final guard atom, whose denotations are the interpretations of the
index terms by \cref{lem:comp-adequacy}.
\end{proof}

\begin{theorem}[Environment validity]
\label{thm:env-validity}
Let $\Env$ be a well-formed definitional environment with $\Sfrag$
declarations, and $G$ a closed shallow goal.  Then:
\begin{enumerate}[leftmargin=2em]
\item (H) holds;
\item $\ML \vDash \ThE$.
\end{enumerate}
\end{theorem}

\begin{proof}
(1) By induction on the position of the definition in $\Env$: the body
of $c \defeq t : \tau_c$ is typed over the strict prefix, for which (H)
holds by induction; \cref{lem:fundamental} with empty context,
discipline and valuation gives $\sem{t}_{} \in \semr{\tau_c}{}$, and
$\cls{c} = \sem{t}_{}$ by the $\delta$-rule.  Lemmas identically via
part (2) of the fundamental lemma.

(2) We check each axiom family of \cref{def:theory}.

\emph{Compilation, pointer and auxiliary axioms}:
\cref{lem:compile-sound}.

\emph{Specification axioms}: \cref{lem:spec-valid}, using (1).

\emph{Premise axioms}: (1) and \cref{lem:coincidence} at the empty
valuation.

\emph{Constructor typing} (datatypes): let $\tup{d}$ be an assignment of
$\tup{y}$ with $\ML \vDash \bigwedge_j \guard{A_{ij}}{y_j} \wedge
\bigwedge_l \F{\varphi_{il}}$ at $\tup{d}$.  By \cref{lem:coincidence},
$d_j \in \semr{A_{ij}}{\rho_0}$ for every $j$ and
$\pvr{\varphi_{il}}{\rho_0} = 1$ for every $l$, where $\rho_0 =
[\tup{y} \mapsto \tup{d}\,]$ (for recursive $j$ the guard is the atom
$\Ihat{I}(y_j, \compf(\erase{\tup{v}_{ij}}))$ and coincidence yields
$(\sem{\tup{v}_{ij}}_{\rho_0}, d_j) \in \semr{I}{}$).  By
\cref{lem:index-reflection}(5), $(\sem{\tup{t}_i}_{\rho_0},
\Ihat{\wcns{c}}_i^{\ML}(\tup{d})) \in \semr{I}{}$, and by
\cref{lem:comp-adequacy} the index arguments of the conclusion atom
denote exactly $\sem{\tup{t}_i}_{\rho_0}$; hence the atom holds.

\emph{Inversion} (datatypes): let $(d, \tup{d}) \in \Ihat{I}^{\ML}$,
i.e.\ $(\tup{d}, d) \in \semr{I}{}$.  \Cref{lem:index-reflection}(1)
yields $i$, $\cls{\tup{e}}$, $\rho_0$; instantiate the existentials of
the $i$-th disjunct with $\cls{\tup{e}}$.  The guards and premise
conjuncts hold by \cref{lem:coincidence} (memberships and truths of the
decomposition); $w \feq \Ihat{\wcns{c}}_i(\tup{y})$ holds by the
interpretation of $\Ihat{\wcns{c}}_i$; and each index equation $z_h
\feq \compf(\erase{t_{ih}})$ holds because $d_h =
(\sem{\tup{t}_i}_{\rho_0})_h$ and \cref{lem:comp-adequacy} identifies
the right-hand side's denotation at the chosen witnesses with
$\sem{t_{ih}}_{\rho_0}$.

\emph{Injectivity and discrimination}: \cref{cor:discrimination}(1) ---
at all elements of $\Dom$, with no guard.

\emph{Predicate clause and case-analysis axioms}: as constructor typing
and inversion one level down, exactly as in \cite{Extraction2026},
using the stage construction of $\pv{J}$, \cref{lem:coincidence} and
\cref{lem:comp-adequacy}.
\end{proof}

\subsection{Main theorem and corollaries}
\label{sec:soundness-main}

\begin{theorem}[Soundness]
\label{thm:soundness}
Let $\Env$ be a well-formed definitional environment with $\Sfrag$
declarations and $G$ a closed shallow proposition over $\Env$.  If
$\ThE \provesfol \F{G}$ then $\pv{G} = 1$: $G$ is true in the
proof-irrelevant classical semantics of $\Env$.
\end{theorem}

\begin{proof}
$\ML \vDash \ThE$ (\cref{thm:env-validity}); by soundness of classical
first-order logic, $\ML \vDash \F{G}$; by \cref{lem:coincidence} at the
empty valuation, $\pv{G} = 1$.
\end{proof}

\begin{corollary}[Consistency]
\label{cor:consistency}
$\ThE$ is consistent, for every well-formed definitional $\Env$ and
goal $G$.  In particular no combination of the emitted axiom families
--- across all declared families, instances and definitions --- can
derive $\bot$.
\end{corollary}

\begin{proof}
$\ML \vDash \ThE$ and $\ML \nvDash \bot$.  No consistency assumption on
the source is needed: a definitional environment always has the model
$\ML$.
\end{proof}

\begin{corollary}[Refutable goals are not provable]
\label{cor:no-refutable}
If $\vdash \pi : \neg G$ for some closed proof over $\Env$, then $\ThE
\nvdash_{\mathrm{FOL}} \F{G}$; more generally $\ThE$ never proves both
$\F{G}$ and $\F{\neg G}$.
\end{corollary}

\begin{proof}
By \cref{lem:fundamental} (via \cref{thm:env-validity}(1)),
$\pv{\neg G} = 1$, so $\pv{G} = 0$ and \cref{thm:soundness} applies
contrapositively; for the second statement, $\ML$ satisfies every
consequence of $\ThE$ and $\F{\neg G} = \F{G} \to \bot$.
\end{proof}

\begin{proposition}[Standardness on purely structural families]
\label{prop:standard}
Call a datatype $I$ \emph{purely structural} if its constructors have no
premises and every non-recursive constructor argument type is an
occurrence of a purely structural datatype declared earlier.  Then:
\begin{enumerate}[leftmargin=2em]
\item For every purely structural $I$ and $(\tup{d}, d) \in \semr{I}{}$,
  $d$ is the class of a closed term built exclusively from constructors
  of purely structural datatypes; and two such closed constructor terms
  have the same class only if they carry the same constructor at the
  root and componentwise equal-class arguments, so each class
  determines its constructor tree over classes uniquely.
\item For $\vect$ over a purely structural element type $A$, the map
  \begin{multline*}
    (a_1, \dots, a_n) \longmapsto\\
    \Bigl(\cls{\csS^n(\csO)},\ \cls{\csvcons(\csS^{n-1}(\csO), \hat a_1,
    \csvcons(\csS^{n-2}(\csO), \hat a_2, \dots, \csvnil \dots))}\Bigr)
  \end{multline*}
  (for representatives $\hat a_i$) is a bijection from the set of finite
  sequences over $\semr{A}{}$ to $\semr{\vect}{}$; in particular
  $\semr{\vect(\csS^n(\csO))}{}$ is in bijection with
  $\semr{A}{}^{\,n}$, and $n \mapsto \cls{\csS^n(\csO)}$ is a bijection
  $\mathbb{N} \to \semr{\nat}{}$.
\end{enumerate}
Consequently, for propositions whose quantifiers range over purely
structural datatypes, $\pv{\cdot}$ is satisfaction in the corresponding
standard structure of index-respecting finite trees.
\end{proposition}

\begin{proof}
(1) By induction on declaration order and, within a datatype, on
$\rank{I}{d}$: \cref{lem:index-reflection}(1) decomposes $d$ as
$\cls{\wcns{c}_i(\tup{e})}$ with non-recursive components in
interpretations of earlier purely structural datatypes (induction on
declaration order) and recursive components of smaller rank (induction
on rank); premises are absent.  Distinctness is
\cref{cor:discrimination}(1) by induction on the terms.

(2) Write $\Phi$ for the displayed map.  \emph{Well-definedness and
injectivity}: the image element is a constructor term; injectivity of
$\Phi$ follows from \cref{cor:discrimination}(1) by induction on
sequence length, together with the injectivity of $n \mapsto
\cls{\csS^n(\csO)}$ (also \cref{cor:discrimination}(1)).  \emph{Image in
$\semr{\vect}{}$}: by induction on $n$ using
\cref{lem:index-reflection}(5) and, for the index computation,
$\sem{\csS\,n}_{[n \mapsto \cls{\csS^{n-1}(\csO)}]} =
\cls{\csS^n(\csO)}$.  \emph{Surjectivity}: let $(\tup{d}, d) \in
\semr{\vect}{}$ and induct on $\rank{\vect}{d}$.  If the decomposition
is by $\csvnil$: $\tup{d} = (\sem{\csO}) = (\cls{\csO})$ and $d =
\cls{\csvnil} = \Phi(\varepsilon)$.  If by $\csvcons$ with components
$\cls{e_n}, \cls{e_a}, \cls{e_w}$: the recursive condition gives
$(\cls{e_n}, \cls{e_w}) \in \semr{\vect}{}$ (the index tuple of the
recursive argument is $(\sem{n}_{\rho_0}) = (\cls{e_n})$), so by the
inner induction $\cls{e_w} = \Phi(a_2, \dots, a_n)$-image with
$\cls{e_n} = \cls{\csS^{n-1}(\csO)}$; moreover $\cls{e_a} \in
\semr{A}{}$, and $\tup{d} = (\sem{\csS\,n}_{\rho_0}) =
(\cls{\csS^{n}(\csO)})$; hence $(\tup{d}, d) = \Phi(\cls{e_a}, a_2,
\dots, a_n)$.  The bijection $\mathbb{N} \cong \semr{\nat}{}$ is the
$m = 0$ instance, as in \cite{Extraction2026}.  The final claim unfolds
\cref{def:interp} clause by clause under the bijections.
\end{proof}

\begin{remark}[Postulated axioms]
\label{rem:axioms}
If $\Env$ additionally contains postulates $c : \varphi$ (or $c :
\sigma$) without bodies, all results relativise as in
\cite{Extraction2026}: \cref{thm:env-validity} goes through
provided each postulate is true in the semantics ($\pv{\varphi} = 1$;
for a postulated constant, $\cls{c} \in \semr{\sigma}{}$ under the
evident extension of $\ML$).  Excluded middle and proof irrelevance are
true; functional extensionality is \emph{not}
(\cref{rem:funext-leak}); a postulated inhabitant of an empty instance
($c : \Fin(\csO)$) is not, and voids the guarantee --- which is the
correct account of the corresponding SMT-datatype hazard
(\cref{ex:fin-empty}): the danger is never the empty instance itself
but an axiom asserting it inhabited.
\end{remark}

\begin{remark}[What is not claimed]
\label{rem:not-claimed}
\Cref{thm:soundness} does not claim $G$ is provable in $\CICi$: the
semantics validates excluded middle, proof irrelevance and (as spent in
the $\tscase{}{}{}{}$ case) uniqueness of identity proofs, so
$\F{G}$ may be provable for $G$ true-but-unprovable.  This matches the
design of the implemented tool, which runs classical provers and relies
on reconstruction.  \Cref{sec:formulas-uip} localises the UIP debt
precisely.
\end{remark}
